% ID: SEL-D0062-OA-E03
% ESTADO: cotejada. Segunda pasada visual y borrador compilado revisados.
% Fuente: Selectividad/sel_2010_matematicas_aplicadas/mat_apli/mat_apli_junio/mat_apli_junio_ra/EXAMEN-ANDALUCÍA-4.pdf
% SHA-256: 4e10e40d60aebdab7570fa77b9ce953d7bf6c9a4c8869a5d3f9b7ca3c907c154
% Página: 1; opción A; ejercicio original 3.

El $41\%$ de quienes se presentan a un examen son varones. Aprueban dicho examen el $70\%$ de los varones presentados y el $60\%$ de las mujeres presentadas.
\begin{enumerate}[label=\alph*)]
\item \textbf{(1 punto)} Calcule la probabilidad de que si una persona escogida al azar ha aprobado, sea mujer.
\item \textbf{(1 punto)} Calcule la probabilidad de que si una persona escogida al azar ha suspendido, sea mujer.
\item \textbf{(0.5 puntos)} Ana dice que si alguien ha aprobado, es más probable que sea mujer que varón; Benito dice que si alguien ha suspendido es más probable que sea mujer que varón. \textquestiondown{}Quién tiene razón?
\end{enumerate}
